\newcommand{\cognitivetable}[5]{%
	\begin{center}
		\renewcommand{\arraystretch}{1.5}
		\begin{tabularx}{\linewidth}{|X|X|}
			\hline 
			\textbf{#1} & \textbf{Osservazioni}
			\\ \hline
			\textbf{1. L'utente sa cosa deve ottenere con l'azione?} (sa cosa deve fare?) & #2 \\ \hline
			\textbf{2. L'utente può individuare lo strumento di interazione?} & #3 \\ \hline
			\textbf{3. L'utente può anche capire che quello è lo strumento giusto per fare ciò che vuol fare?} & #4 \\ \hline
			\textbf{4. Dopo che l'azione è eseguita, l'utente capisce la risposta che ottiene?} & #5 \\ \hline
		\end{tabularx}
	\end{center}
}
\newcommand{\tabellaProblemi}[6]{
	\noindent
	\begin{tabularx}{\textwidth}{|X|}
		\hline
		\begin{minipage}[t]{\linewidth}\raggedright
			\vspace{4pt}
			\textbf{Problema N\textsuperscript{o}:} #1 \hfill \textbf{Grado di Severità:} #5\\[4pt]
			\textbf{Azione:} #2\\[6pt]
			\textbf{Domanda:} \par
			#3\\[6pt]
			\textbf{Descrizione:} \par
			#4
			\vspace{6pt}
		\end{minipage} \\ \hline
		
		\begin{minipage}[t]{\linewidth}\centering
			\vspace{4pt}
			\textbf{Screenshot:}
			\vspace{0.6em}
			
			#6
			\vspace{6pt}
		\end{minipage} \\ \hline
	\end{tabularx}
	\vspace{1em}
}

\definecolor{negativeRed}{RGB}{180, 40, 40}
\definecolor{orangino}{RGB}{205, 127, 50}

\section{Walkthrough Cognitivo}
Abbiamo deciso di eseguire un walkthrough cognitivo focalizzato principalmente sulle azioni eseguibili da un cliente, in quanto unico individuo dal quale non ci si aspetta alcuna preparazione verso il nostro sistema: ogni nuovo cliente si dovrà affacciare ad un nuovo sistema, il quale deve permettere ad ognuno di adattarsi velocemente all'interfaccia mostrata.

Le azioni che siamo andati ad analizzare sono le seguenti:
\begin{itemize}
	\item Fare un bonifico;
	\item Reimpostare la password.
\end{itemize}
Altre azioni non sono state prese in considerazione, in quanto già considerate nel capitolo 3 oppure incluse/simili alle azioni qua mostrate (esempio di Modifica Profilo, che presenta le stesse caratteristiche di Reimpostare Password, controlli a parte).

\subsection{Prerequisiti}
La fase preliminare consiste nella definizione degli elementi principali necessari per l'analisi del walkthrough:
\begin{itemize}
	\item \textbf{Sistema utilizzato}: abbiamo deciso di utilizzare il sistema nella forma antecedente alla riprogettazione.
	\item \textbf{Indicazione sugli utenti}: i valutatori impersonano un utente.
\end{itemize}
Viene riportato il documento della lista azioni reazioni utilizzato:\\
\begin{center}
	\includegraphics[width=0.50\linewidth]{listaazionireazioni}
\end{center}
Viene riportata la tabella utilizzata per il documento di valutazione:\\

\cognitivetable{Nome dell'Azione}{}{}{}{}

\subsection{Fase di valutazione}
Di seguito, le tabelle per Reimpostare password.

\cognitivetable{%
	Andare nella sezione del profilo 
}{
	Sì. L'utente sa \textcolor{negativeRed}{(forse)} che deve accedere alla sezione dedicata ai suoi dati per trovare le funzionalità relativa alla modifica della password.%
}{%
	Sì, il pulsante "Il Mio Profilo" è evidente, è uno dei primi elementi che compaiono a video quando accedo alla dashboard.
}{%
	Sì, un pulsante giallo "Il Mio Profilo" è abbastanza autoesplicativo per poter comprendere che rimanda ai dettagli del proprio profilo.
}{%
	Sì. Il sistema mostra la pagina relativa alle informazioni dell'utente, con le relative funzionalità associate.%
}

\cognitivetable{%
	Andare nella sezione modifica password
}{
	{Sì. Cliccando su modifica password sa che andrà a modificare la password.}%
}{%
	Sì, il pulsante è molto visibile.
}{%
	\textcolor{negativeRed}{No, il pulsante è ripetuto due volte e confondibile con un altro pulsante chiamato "Sicurezza".}
}{%
	Sì. Il sistema mostra la pagina relativa alla modifica della password, con il form dedicato per l'inserimento dei dati (vecchia password, nuova password, conferma nuova password).%
}

\cognitivetable{%
	Inserire la password
}{
	Sì. L'utente sa che sta inserendo i dati (vecchia password, nuova password, conferma nuova password) sta andando a cambiare la sua password (titolo della sezione in alto).
}{%
	Sì, il form è al centro della pagina.
}{%
	Sì, il titolo del form è "Cambio password".
}{%
	Sì, il sistema comunica eventuali errori o conferma dell'avvenuto cambiamento all'utente.%
}

\cognitivetable{%
	Confermare l'operazione
}{
	\textcolor{negativeRed}{No, il sistema non chiede conferma all'utente.}
}{%
	\textcolor{negativeRed}{No, il sistema non chiede conferma all'utente.}
}{%
	\textcolor{negativeRed}{No, il sistema non chiede conferma all'utente.}
}{%
	\textcolor{negativeRed}{No, il sistema non chiede conferma all'utente.}
}

Di seguito, le tabelle per Fare un bonifico.

\cognitivetable{%
	Passare alla sezione per fare un bonifico.
}{
	Sì è chiaro.
}{%
	Sì il bottone "Bonifio" è chiaramente visibile %
}{%
	Sì in quanto non ci sono ambiguità%
}{%
	Sì. Il sistema mostra la pagina per eseguire il bonifico%
}

\noindent

\cognitivetable{%
	Inserire i dati del beneficiario
}{
	Sì è chiaro.
}{%
	Sì il forum per l'input è evidente%
}{%
	Sì in quanto non ci sono ambiguità%
}{%
	\textcolor{orangino}{Sì e no. Il testo che conferma la correttezza dei dati inseriti è illeggibile oppure non si nota}%
}

\clearpage
\noindent

\cognitivetable{%
	Inserire l’importo e causale
}{
	Sì è chiaro.
}{%
	Sì il forum per l'input è evidente%
}{%
	Sì in quanto non ci sono ambiguità%
}{%
	Sì appare un elemento visivo che fa capire all'utente che ha svolto tutte le operazioni di inserimento correttamente%
}

\noindent

\cognitivetable{%
	Confermare o annullare il bonifico
}{
	Sì è chiaro.
}{%
	Sì il bottone per confermare diventa abilitato. Per quanto riguarda l'annullamento anche questo bottone è chiaramente visibile vicino al bottone di conferma con una certa spaziatura fra i due.%
}{%
	Sì in quanto non ci sono ambiguità%
}{%
	Sì appare una schermata che mostra che il bonifico è andato a buon fine%
}

\subsection{Fase di Debriefing}
Abbiamo stilato un documento unificato dei problemi individuati, qui riportati:

\tabellaProblemi{
	1 
}{
	Reimpostare Password - Confermare l'operazione (cambio password). 
}{
	Tutte.
}{
	Quando l'utente clicca su cambia password, il sistema non chiede conferma: se i dati sono corretti la cambia subito, se non sono corretti dà errore. 
}{
	4 
}{
	\textit{Nessuno screenshot allegato} 
}

\tabellaProblemi{
	2 
}{
	Reimpostare Password - Andare nella sezione modifica password. 
}{
	L'utente può anche capire che quello è lo strumento giusto per fare ciò che vuol fare?
}{
	Presenti due tasti identici "Cambia Password", più dei tasti per la gestione delle domande di sicurezza confondibili (già visto nella valutazione euristica).
}{
	3 
}{
	\includegraphics[width=0.75\linewidth]{problema_14} 
}

\tabellaProblemi{
	3 
}{
	Reimpostare Password - Andare nella sezione del profilo. 
}{
	L'utente sa cosa deve ottenere con l'azione? (sa cosa deve fare?)
}{
	Il tasto per andare nel profilo è chiaro, è da capire se l'utente può giungere facilmente alla conclusione "nel mio profilo potrò cambiare password".
}{
	1
}{
	\textit{Nessuno screenshot allegato} 
}

\tabellaProblemi{
	4
}{
	Fare un bonifico - Inserire i dati del beneficiario
}{
	4
}{
	Dopo aver inserito i dati del destinatario l'utente ha un feedback praticamente inesistente 
}{
	2 
}{
	\includegraphics[width=0.8\textwidth]{problema_1_foto_1}
	\includegraphics[width=0.8\textwidth]{problema_1_foto_2}
}